% ==================================================
% PREAMBLE
% ==================================================
\documentclass[12pt]{article}

% -------- Packages --------
\usepackage[utf8]{inputenc}
\usepackage{natbib}
\usepackage{setspace}
\usepackage[english]{babel}
\usepackage{amsmath}
\usepackage{graphicx}
\usepackage[colorlinks=true, allcolors=blue]{hyperref}
\hypersetup{
    colorlinks=true,
    linkcolor=black,
    filecolor=magenta,      
    urlcolor=cyan,
    pdfpagemode=FullScreen
}
\usepackage{wasysym}
\usepackage{mathabx}
\usepackage{enumitem}
\usepackage{comment}
\usepackage[table]{xcolor}
\usepackage{geometry}
\geometry{legalpaper,margin=0.7in}    % margin left and right is 1inch
% \usepackage[legalpaper,top=1cm,bottom=2cm,left=0.5cm,right=0.5cm,marginparwidth=1.75cm]{geometry}
\usepackage{wrapfig}
\usepackage{caption}
\usepackage{subcaption}
\usepackage{multicol}
\usepackage{fancyhdr}
\usepackage{datetime2}   % better date formatting
\usepackage{tabularx}
\usepackage{multirow}
\usepackage{array}
\usepackage{makecell}
\usepackage[most]{tcolorbox}

% -------- Custome Commands --------

\pagestyle{fancy}
\fancyhf{}  % clear default header/footer
\fancyhead[L]{COHA Cheatsheet} % header content
\fancyhead[R]{Last edited: \today} % header content
\fancyfoot[R]{\thepage} % page number in footer

\DTMsetdatestyle{ddmmyyyy}

\newcounter{row}
\newcolumntype{N}{>{\stepcounter{row}\arabic{row}}c}

\setlength{\parindent}{0pt}

%%%%%%%%%%%%%%%%%%%%%%%%%%%%%%%%%%%%%%%%%%%%%%%%%%%%
% HEADER + TOC
%%%%%%%%%%%%%%%%%%%%%%%%%%%%%%%%%%%%%%%%%%%%%%%%%%%%

\begin{document}
{
  \hypersetup{linkcolor=black}
  \tableofcontents
}
\clearpage

%%%%%%%%%%%%%%%%%%%%%%%%%%%%%%%%%%%%%%%%%%%%%%%%%%%%
% CONTENT
%%%%%%%%%%%%%%%%%%%%%%%%%%%%%%%%%%%%%%%%%%%%%%%%%%%%

%%%%%%%%%%%%%%%%%%%%%%%%%%%%%%%%%%%%%%%%%%%%%%%%%%%%
\section{Anleitungen}

\subsection{Bagels}
\begin{tabularx}{\textwidth}{|l|X|}
    \hline
    \rowcolor{gray!25}
    \multicolumn{2}{|l||}{Schinken-Käse}\\
    \hline
    Frischkäse & auf Deckel und Boden \\
    Zupfsalat & \\
    Schinken & 2 Blätter \\
    Käse & 1 Scheibe \\
    Gurke & ca. 3 Scheiben\\
    Tomaten & ca. 3 Scheiben \\ 
    Salz \& Pfeffer & \\ \hline
\end{tabularx}

\vspace{1em}

\begin{tabularx}{\textwidth}{|l|X|}
    \hline
    \rowcolor{gray!25}
    \multicolumn{2}{|l|}{Käse-Gurke} \\
    \hline
    Frischkäse & auf Deckel und Boden \\
    Zupfsalat & \\
    Gurke & 5-6 Scheiben \\
    Käse & 2 Scheiben \\
    Salz \& Pfeffer & \\ \hline
\end{tabularx}

\vspace{1em}

\begin{tabularx}{\textwidth}{|l|X|}
    \hline
    \rowcolor{gray!25}
    \multicolumn{2}{|l|}{Lachs} \\
    \hline
    Frischkäse & auf Deckel und Boden \\
    Senf & Teelöffel auf Deckel od. Boden verstreichen \\
    Gurke & ca. 5 Scheiben \\
    Lachs & ca. 2 Stk. \\
    Rucola & \\
    Salz \& Pfeffer & \\ \hline
\end{tabularx}

\vspace{1em}

\begin{tabularx}{\textwidth}{|l|X|}
    \hline
    \rowcolor{gray!25}
    \multicolumn{2}{|l|}{Tomaten - Mozzarella} \\
    \hline
    Pesto & auf Deckel und Boden \\
    Tomaten & 4-5 Scheiben \\
    Mozzarella & 2 Scheiben \\
    Rucola & \\
    Salz \& Pfeffer & \\ \hline
\end{tabularx}


\clearpage
%%%%%%%%%%%%%%%%%%%%%%%%%%%%%%%%%%%%%%%%%%%%%%%%%%%%
\subsection{Bagels}

\begin{minipage}[t]{0.48\textwidth}
    \begin{tcolorbox}[title = Schinken - Käse, left=0.5pt, arc=0pt]
        \begin{itemize}[itemsep=1pt, topsep=0pt]
            \item[-] Frischkäse
            \item[-] Schinken
            \item[-] Käse
            \item[-] Zupfsalat
            \item[-] Gurke
            \item[-] Tomate
            \item[-] Salz \& Pfeffer
        \end{itemize}
    \end{tcolorbox}
\end{minipage}
\hfill
\begin{minipage}[t]{0.48\textwidth}
    \begin{tcolorbox}[title=Gurke - Käse, left=0.5pt, arc=0pt]
        \begin{itemize}[itemsep=1pt, topsep=0pt]
            \item[-] Frischkäse
            \item[-] Käse
            \item[-] Zupfsalat
            \item[-] Gurke
            \item[-] Salz \& Pfeffer
        \end{itemize}
    \end{tcolorbox}
\end{minipage}

\vspace{1em}

\begin{minipage}[t]{0.48\textwidth}
    \begin{tcolorbox}[title = Lachs, left=0.5pt, arc=0pt]
        \begin{itemize}[itemsep=1pt, topsep=0pt]
            \item[-] Frischkäse
            \item[-] Senf
            \item[-] Lachs
            \item[-] Rucola
            \item[-] Gurke
            \item[-] Salz \& Pfeffer
        \end{itemize}
    \end{tcolorbox}
\end{minipage}
\hfill
\begin{minipage}[t]{0.48\textwidth}
    \begin{tcolorbox}[title= Tomaten - Mozzarella, left=0.5pt, arc=0pt]
        \begin{itemize}[itemsep=1pt, topsep=0pt]
            \item[-] Pesto
            \item[-] Mozzarella
            \item[-] Rucola
            \item[-] Tomaten
            \item[-] Salz \& Pfeffer
        \end{itemize}
    \end{tcolorbox}
\end{minipage}

\vspace{1em}

\begin{minipage}[t]{0.48\textwidth}
    \begin{tcolorbox}[title = Creamcheese, left=0.5pt, arc=0pt]
        \begin{itemize}[itemsep=1pt, topsep=0pt]
            \item[-] Frischkäse
            \item[-] Rucola
            \item[-] Gurke
            \item[-] Käse 
            \item[-] Salz \& Pfeffer
        \end{itemize}
    \end{tcolorbox}
\end{minipage}


\clearpage
%%%%%%%%%%%%%%%%%%%%%%%%%%%%%%%%%%%%%%%%%%%%%%%%%%%%
\subsection{Bagels}

\begin{multicols}{2}

\textbf{Schinken - Käse }
\begin{itemize}[itemsep=1pt, topsep=5pt]
\item Frischkäse
\item Schinken
\item Käse
\item Zupfsalat
\item Gurke
\item Tomate
\item Salz \& Pfeffer
\end{itemize}
\columnbreak

\textbf{Gurke - Käse }
\begin{itemize}[itemsep=1pt, topsep=5pt]
\item Frischkäse
\item Käse
\item Zupfsalat
\item Gurke
\item Salz \& Pfeffer
\end{itemize}

\end{multicols}
\begin{multicols}{2}

\textbf{Lachs}
\begin{itemize}[itemsep=1pt, topsep=5pt]
\item Frischkäse
\item Senf
\item Lachs
\item Rucola
\item Gurke
\item Salz \& Pfeffer
\end{itemize}
\columnbreak

\textbf{Tomate - Mozzarella}
\begin{itemize}[itemsep=1pt, topsep=5pt]
\item Pesto
\item Tomate
\item Mozzarella
\item Rucola
\item Salz \& Pfeffer
\end{itemize}
\end{multicols}


\clearpage
%%%%%%%%%%%%%%%%%%%%%%%%%%%%%%%%%%%%%%%%%%%%%%%%%%%%
\subsection{Kaffee}
\begin{tabularx}{\textwidth}{|X|c|c|c|c|l|X|}
    \hline
    Typ &
    Espresso &
    \makecell{Wasser\\(y/n)} &
    \makecell{Milch \\ (y/n)} &
    \makecell{Schaum \\ (y/n)} &
    Tasse &
    Zusatz-infos\\
    \hline
    Espresso    & 1     & n & n & n & klein     & ='Kurzer' \\\arrayrulecolor{gray}\hdashline\arrayrulecolor{black}
    Doppio      & 2     & n & n & n & mittel    & \\\arrayrulecolor{gray}\hdashline\arrayrulecolor{black}
    Macchiato   & 1     & n & n & y & klein & \\\arrayrulecolor{gray}\hdashline\arrayrulecolor{black}
    Verlängerter & 1    & y & n & n & mittel    & \\\arrayrulecolor{gray}\hdashline\arrayrulecolor{black}
    Milchkaffe  & 1     & y & y & n & mittel    & \\\arrayrulecolor{gray}\hdashline\arrayrulecolor{black}
    Cappuccino  & 1     & n & y & y & groß      & \\\arrayrulecolor{gray}\hdashline\arrayrulecolor{black}
    Melange     & 1     & y & y & y & groß      & \\\arrayrulecolor{gray}\hdashline\arrayrulecolor{black}
    Flat White  & 2     & n & y & y & groß      & =cappu+extra shot \\\arrayrulecolor{gray}\hdashline\arrayrulecolor{black}
    Cafe Latte  & 1     & n & y & n & Glas mittel & mit Schaum nicht so genau \\\arrayrulecolor{gray}\hdashline\arrayrulecolor{black}
    großer Brauner & 2  & n & y & n & mittel    & normal Milch extra \\\arrayrulecolor{gray}\hdashline\arrayrulecolor{black}
    kleiner Brauner & 1 & n & y & n & klein     & normal Milch extra \\\arrayrulecolor{gray}\hdashline\arrayrulecolor{black}
    Eiskaffe    & 1     & y & n & n & Glas groß & +2Kgl Vanille Eis +Schlag \\\arrayrulecolor{gray}\hdashline\arrayrulecolor{black}
    Affogato    & 1     & n & n & n & mittel    & +1Kgl Vanille Eis \\
    \hline
\end{tabularx}

\clearpage
%%%%%%%%%%%%%%%%%%%%%%%%%%%%%%%%%%%%%%%%%%%%%%%%%%%%
\section{Checklisten}

\subsection{Frühdienst}
\begin{enumerate}
    \item Im Kalender schauen ob es Reservierungen gibt
    \item Backofen, Geschirrspüler, Eisvitrine einschalten
    \item Kuchenvitrine herrichten
    \item Bagels herrichten
    \item Croissants aufbacken
    \item Eisvitrine herrichten (Eis von hinten holen)
    \item Gastgarten herrichten
    \begin{itemize}
        \item Tische abwischen
        \item Schirme aufmachen
        \item Tischnummern, Karten und Zuckergläser(evtl. noch auffüllen) platzieren
    \end{itemize}
    \item gegebenenfalls Licht einschalten
    \item Kaffeemaschine aktivieren
\end{enumerate}

\subsection{Abenddienst}
\begin{enumerate}
    \item restliches Personalessen eintragen
    \item Abrechnung (erst ab 18Uhr und keine offenen Tische mehr)
    \item Gastgarten wieder abbauen (Schirme abspannen; Tischnummern, Karten, Zuckergläser, Aschenbecher weg)
    \item Getränkeladen auffüllen
    \item Leergut wegbringen
    \item Backofen ausschalten und reinigen
    \item Kuchenvitrine ausräumen und säubern
    \item Croissants und Bagels für den nächsten Tag einpacken und in Kühlladen geben.
    \begin{itemize}
        \item[] Außer am Sonntag: da am nächsten Tag(Montag) geschlossen ist, können die übrig gebliebenen Croissants mitgenommen werden. Dann auf Personal/Bruch buchen \textcolor{red}{$\longrightarrow$ TODO!}
    \end{itemize}
    \item Schank Putzen
    \item Milchsystem reinigen
    \item Kaffeemaschine reinigen
    \item Müll weg bringen und frische Müllsäcke in Mülleimer geben
    \item Oberflächen reinigen
    \item Eisvitrine ausschalten und Eis wieder nach hinten ins Lager bringen oder in kleinen TK geben, wenn Platz hat
    \item Geschirrspüler ausschalten/abpumpen/Siebe reinigen
    \item Licht ausschalten
    \item Klos zu sperren und alle Türen zu sperren
\end{enumerate}


\clearpage
%%%%%%%%%%%%%%%%%%%%%%%%%%%%%%%%%%%%%%%%%%%%%%%%%%%%
\section{Checklisten}

\addcontentsline{toc}{subsection}{Frühdienst}
\begin{tabularx}{\textwidth}{|N|X r|}
    \hline
    \rowcolor{gray!25}
    \multicolumn{2}{|l}{\large{\textbf{Frühdienst}}\rule{0pt}{2.5ex}} & Info \\
    \hline
    & Reservierungen checken & im Kalender am iPad\\
    & Backofen, Geschirrspüler, Eisvitrine einschalten & brauchen alle Zeit bis einsatzbereit sind \\
    & Kuchenvitrine herrichten & \\
    & Bagels herrichten & zuerst die vom Vortag verwenden \\
    & Croissants aufbacken & zuerst die vom Vortag verwenden \\
    & Eisvitrine herrichten & \\
    & Gastgarten herrichten & \makecell[l]{Tische reinigen, Schirme aufspannen, \\ Tischnummern + Karten + Zuckergläser aufstellen }\\
    & gegebenenfalls Licht einschalten &\\
    & Kaffeemaschine aktivieren &\\
    & Klos aufsperren & wenn nicht schon von Reinigungspersonal aufgesperrt\\
    \hline
\end{tabularx}

\vspace{1em}

\addcontentsline{toc}{subsection}{Abbenddienst}
\setcounter{row}{0}
\begin{tabularx}{\textwidth}{|N|X r|}
    \hline
    \rowcolor{gray!25}
    \multicolumn{2}{|l}{\large{\textbf{Abenddienst}}\rule{0pt}{2.5ex}} & Zusatzinfo \\
    \hline
    & restliches Personalessen eintragen &\\
    & Abrechnung & erst ab 18Uhr und keine offenen Tische mehr\\
    & Gastgarten abbauen &\\
    & Getränkeladen auffüllen & \\
    & Leergut wegbringen & \\
    & Backofen ausschalten und reinigen & \\
    & Kuchenvitrine ausräumen und säubern & \\
    & Croissants und Bagels für den nächsten Tag einpacken und in Kühlladen geben. & nicht am Sonntag! \\
    & Schank putzen &\\
    & Milchsystem reinigen &\\
    & Kaffeemaschine reinigen &\\
    & Müll weg bringen, frische Müllsäcke in Eimer geben&\\
    & Oberflächen reinigen &\\
    & Eisvitrine ausschalten, Eis wieder in Lager/kleinen TK geben &\\
    & Geschirrspüler ausschalten/abpumpen/Siebe reinigen & Türe offen lassen\\
    & Licht ausschalten &\\
    & Umsatzabgabe in Briefkasten werfen &\\
    & Alle Türen + Klos zu sperren &\\
    \hline
\end{tabularx}

\end{document}